\begin{definition}[Правильная раскраска и хроматический многочлен]
Правильная раскраска в \(k\) цветов — функция \(\varphi:V\to\{1,\ldots,k\}\),
для которой \(\varphi(u)\ne\varphi(v)\) при \(uv\in E\).
Разрешены неиспользованные цвета; цвета различимы.
Число таких раскрасок обозначается \(\chi_G(k)\).
Хроматическое число \(\chi(G)\) — минимальное число цветов правильной раскраски.
\end{definition}
\begin{lemma}[Удаление и стягивание]
Для ребра \(e=uv\) простого графа
\[
 \boxed{\chi_G(k)=\chi_{G-e}(k)-\chi_{G/e}(k).}
\]
После стягивания кратные рёбра можно заменить одиночными.
\end{lemma}
\begin{proof}
Разобьём раскраски \(G-e\) на два класса.
При различных цветах \(u,v\) это ровно раскраски \(G\).
При одинаковых цветах \(u,v\) отождествляем их: получается биекция
с раскрасками \(G/e\). Кратные рёбра задают то же самое неравенство цветов,
поэтому их число на раскраски не влияет. Складывая мощности, получаем формулу.
\end{proof}
\begin{graphillustration}{Лекционный пример удаления и стягивания одного ребра квадрата. Слева сохраняются четыре вершины, справа концы $u,v$ отождествлены.}{fig:18-delete-contract}
\begin{tikzpicture}[font=\small]
\node[gi vertex] (u) at (0,1.6) {$u$};
\node[gi vertex] (v) at (0,0) {$v$};
\node[gi vertex] (a) at (1.6,0) {$a$};
\node[gi vertex] (b) at (1.6,1.6) {$b$};
\draw[gi edge] (v) -- (a);
\draw[gi edge] (a) -- (b);
\draw[gi edge] (b) -- (u);
\draw[gi removed] (u) -- (v);
\node[gi note] at (0.8,-0.8) {$G-e=P_4$};
\node[gi note] at (3,0.8) {$\longleftarrow$};
\begin{scope}[shift={(4.2,0)}]
\node[gi vertex] (u) at (0,1.6) {$u$};
\node[gi vertex] (v) at (0,0) {$v$};
\node[gi vertex] (a) at (1.6,0) {$a$};
\node[gi vertex] (b) at (1.6,1.6) {$b$};
\draw[gi emphasis] (u) -- (v);
\draw[gi edge] (v) -- (a);
\draw[gi edge] (a) -- (b);
\draw[gi edge] (b) -- (u);
\node[gi note] at (0.8,-0.8) {$G=C_4$, $e=uv$};
\end{scope}
\node[gi note] at (7.2,0.8) {$\longrightarrow$};
\begin{scope}[shift={(8.5,0)}]
\node[gi selected] (w) at (0,0.8) {$u=v$};
\node[gi vertex] (a) at (1.6,0) {$a$};
\node[gi vertex] (b) at (1.6,1.6) {$b$};
\draw[gi edge] (w) -- (a);
\draw[gi edge] (a) -- (b);
\draw[gi edge] (b) -- (w);
\node[gi note] at (0.8,-0.8) {$G/e=C_3$};
\end{scope}
\node[gi note] at (5,-1.8) {$\chi_{C_4}=\chi_{P_4}-\chi_{C_3}$};
\end{tikzpicture}
\end{graphillustration}


\begin{graphillustration}{Два класса раскрасок $G-e$: показано по одному примеру. Цифры — цвета; равенство цветов концов удалённого ребра соответствует стягиванию.}{fig:18-color-classes}
\begin{tikzpicture}[font=\small]
\begin{scope}[shift={(0,0)}]
\node[gi selected,label=left:$u$] (u) at (0,1.6) {1};
\node[gi selected,label=left:$v$] (v) at (0,0) {2};
\node[gi vertex,label=right:$a$] (a) at (1.6,0) {1};
\node[gi vertex,label=right:$b$] (b) at (1.6,1.6) {2};
\draw[gi edge] (v) -- (a);
\draw[gi edge] (a) -- (b);
\draw[gi edge] (b) -- (u);
\node[gi note] at (0.8,-0.9) {$\varphi(u)\ne\varphi(v)$: раскраски $G$};
\end{scope}
\begin{scope}[shift={(6,0)}]
\node[gi selected,label=left:$u$] (u) at (0,1.6) {1};
\node[gi selected,label=left:$v$] (v) at (0,0) {1};
\node[gi vertex,label=right:$a$] (a) at (1.6,0) {3};
\node[gi vertex,label=right:$b$] (b) at (1.6,1.6) {2};
\draw[gi edge] (v) -- (a);
\draw[gi edge] (a) -- (b);
\draw[gi edge] (b) -- (u);
\node[gi note] at (0.8,-0.9) {$\varphi(u)=\varphi(v)$: раскраски $G/e$};
\end{scope}
\end{tikzpicture}
\end{graphillustration}


\begin{cor}[Добавление и стягивание]
Для несмежных \(u,v\)
\[
 \chi_G(k)=\chi_{G+uv}(k)+\chi_{G/(u=v)}(k).
\]
\end{cor}
\begin{proof}
Снова разделяем раскраски по равенству цветов \(u,v\).
Разные цвета соответствуют \(G+uv\), одинаковые — отождествлению \(u,v\).
\end{proof}
\begin{theorem}[Полиномиальность]
Для графа без петель \(\chi_G(k)\) является унитарным многочленом
степени \(n\) с целыми коэффициентами. Для пустого множества вершин \(\chi_G=1\).
\end{theorem}
\begin{proof}
Индукция по \(n+m\); кратные рёбра предварительно удаляем.
Для графа без рёбер вершины красятся независимо, \(\chi_G(k)=k^n\).
При наличии ребра используем удаление–стягивание: первый многочлен
по индукции унитарен степени \(n\), второй — степени \(n-1\).
Их разность унитарна степени \(n\) и имеет целые коэффициенты.
Полученный многочлен совпадает с числом раскрасок для всех натуральных \(k\).
Он единственен: ненулевой многочлен не может иметь бесконечно много корней.
\end{proof}
\begin{lemma}[Произведение по компонентам]
Если компоненты графа — \(G_1,\ldots,G_c\), то
\(\chi_G(k)=\prod_{j=1}^c\chi_{G_j}(k)\).
\end{lemma}
\begin{proof}
Выбор раскраски каждой компоненты независим, поскольку рёбер между
компонентами нет. Правило произведения даёт равенство для натуральных \(k\),
а значит, и тождество многочленов.
\end{proof}
\begin{example}
\(\chi_{K_n}(k)=k(k-1)\cdots(k-n+1)\): последовательно выбираем различные
цвета для всех вершин. Для цикла \(C_3\) рекурсия даёт
\(\chi_{C_3}=k(k-1)^2-k(k-1)=k(k-1)(k-2)\).
Лекционный пример \(C_4\): удаление ребра даёт \(P_4\), стягивание — \(C_3\), поэтому
\[
\chi_{C_4}(t)=t(t-1)^3-t(t-1)(t-2)=t(t-1)(t^2-3t+3).
\]
При прямом подсчёте фиксируем цвет первой вершины. Если противоположная
имеет тот же цвет, получается \(t(t-1)^2\) раскрасок; если другой —
\(t(t-1)(t-2)^2\). Сумма даёт ту же формулу.
Петля запрещает правильные раскраски: \(\chi_G\equiv0\).
\end{example}
